\documentclass[12pt]{article}
\usepackage[margin=0.58in,top=0.62in,bottom=0.48in]{geometry}
\usepackage{lmodern}
\usepackage{microtype}
\usepackage{fancyhdr}
\usepackage{titlesec}
\usepackage{enumitem}
\usepackage{lastpage}
\usepackage[hidelinks]{hyperref}

\setlength{\parindent}{0pt}
\setlength{\parskip}{0.18em}
\setlength{\headheight}{26pt}
\raggedbottom
\titleformat{\section}{\normalsize\bfseries\scshape}{}{0pt}{}
\titleformat{\subsection}{\normalsize\bfseries}{}{0pt}{}
\titlespacing*{\section}{0pt}{0.35em}{0.12em}
\titlespacing*{\subsection}{0pt}{0.25em}{0.08em}
\pagestyle{fancy}
\fancyhf{}
\fancyhead[C]{\footnotesize Lit Example Reader \quad Assignment ID: U1L10LE \quad Page \thepage\ of \pageref{LastPage}}
\renewcommand{\headrulewidth}{0.5pt}

\newenvironment{playpassage}{\begin{quote}\footnotesize\setlength{\parskip}{0.08em}\setlength{\parindent}{0pt}}{\end{quote}}
\newcommand{\speaker}[1]{\textbf{#1.}}

\begin{document}

\begin{center}
{\LARGE\bfseries Week 10 Lit Example Reader}\par\vspace{0.02em}
{\large\scshape Julius Caesar: Proof That Misses the Point}\par\vspace{0.06em}
\rule{0.78\textwidth}{0.5pt}
\end{center}

\textbf{Purpose:} Julius Caesar passages for Week 10: test whether a speech proves the exact point at issue or only a nearby claim.

\textbf{What to watch for:} Strong evidence, moving examples, and public pressure that may answer a different question from the one under debate.

\section*{Before the Passages: The Week 10 Logic Tool}

\begin{itemize}[leftmargin=1.3em,itemsep=0.05em,topsep=0.08em,parsep=0.03em]
\item \textbf{Point at issue:} the exact question the speaker must answer.
\item \textbf{Nearby claim:} something related and maybe true, but not the same as the conclusion.
\item \textbf{Repair:} add evidence that answers the issue, or narrow the conclusion to match what was actually proved.
\end{itemize}

\section*{Passage 1: Act 3, Scene 2 --- Brutus Explains the Killing}

\textbf{Context:} Brutus must justify Caesar's assassination to the Roman people. Watch what he proves, and what still needs proof.

\begin{playpassage}
\speaker{Brutus} Not that I loved Caesar less, but that I loved Rome more.\\
Had you rather Caesar were living and die all slaves, than that Caesar\\
were dead, to live all free men? As Caesar loved me, I weep for him;\\
as he was fortunate, I rejoice at it; as he was valiant, I honour him:\\
but, as he was ambitious, I slew him. There is tears for his love; joy\\
for his fortune; honour for his valour; and death for his ambition.\\
Who is here so base that would be a bondman? If any, speak; for him\\
have I offended.
\end{playpassage}

\subsection*{Reader Notes}

This is your assisted example. Brutus gives the crowd a clear public question: was Caesar's death justified? He offers a nearby claim: Caesar was ambitious, and ambition threatens liberty. Week 10 asks whether the evidence proves the exact conclusion---that killing Caesar was necessary and just---or whether it proves only a related concern about ambition and freedom. Notice the slave/free alternative: it can make the issue feel settled before all parts of the case have been proved.

\textbf{Reading task:} Underline Brutus's main conclusion. Box the word \textit{ambitious}. In the margin, write the exact point at issue: what must Brutus prove to justify the killing?

\newpage

\section*{Passage 2: Act 3, Scene 2 --- Antony Answers}

\textbf{Context:} Antony replies after Brutus. Watch whether Antony proves the conspirators wrong, proves Caesar was not ambitious, or proves a nearby emotional point.

\begin{playpassage}
\speaker{Antony} He hath brought many captives home to Rome\\
Whose ransoms did the general coffers fill:\\
Did this in Caesar seem ambitious?\\
When that the poor have cried, Caesar hath wept:\\
Ambition should be made of sterner stuff:\\
Yet Brutus says he was ambitious;\\
And Brutus is an honourable man.\\
You all did see that on the Lupercal\\
I thrice presented him a kingly crown,\\
Which he did thrice refuse: was this ambition?
\end{playpassage}

\subsection*{Reader Notes}

This passage shifts more work to you. Antony gives evidence that may answer part of Brutus's charge: ransom money for Rome, pity for the poor, and refusing a crown. Mark what each piece of evidence actually proves. Then decide what larger conclusion the crowd may be moving toward, and whether Antony's proof fully reaches it.

\textbf{Reading task:} Number Antony's three pieces of evidence. Next to each, write the smaller claim it supports. Save the final prompt-match judgment for the worksheet.

\section*{How to Use These Passages on the Worksheet}

Do not summarize the funeral speeches. Use Week 10's question: what exact point is at issue, and does the offered proof match it?

\end{document}
